% !TEX program = xelatex
% ============================================================
% 配方：概念示意图（机理/几何/物理关系类，对标"轨道六根数几何定义"类图）
% 编译：xelatex tikz_schematic.tex → 产出 PDF → 正文 \includegraphics 引入
% 用途：题目机理本身需要"画出来才讲得清"时（坐标系/几何关系/作用链路），
%       替代 matplotlib 硬画的扁平示意图。
% 本模板演示可复用词汇：坐标轴 / 平面 / 矢量 / 角度弧 / 标注引线 / 参数框。
% ============================================================
\documentclass[border=8pt]{standalone}
\usepackage{ctex}
\usepackage{newtxmath}
\usepackage{tikz}
\usetikzlibrary{arrows.meta,angles,quotes,positioning,calc}

\begin{document}
\begin{tikzpicture}[
  font=\small,
  axis/.style  = {-{Stealth[length=2.2mm]}, line width=0.8pt},
  vec/.style   = {-{Stealth[length=2.2mm]}, line width=1.1pt, blue!60!black},
  aux/.style   = {dashed, black!55, line width=0.5pt},
  lbl/.style   = {font=\footnotesize},
  callout/.style={font=\footnotesize, draw=gray!60, fill=white, rounded corners=1mm,
                  inner sep=2pt, align=center}
]

% ---- 参考平面（如：基准面/可行域，按题目改） ----
\fill[blue!4] (0,0) -- (4.2,0.4) -- (5.2,2.2) -- (1.0,1.8) -- cycle;
\draw[black!45, line width=0.5pt] (0,0) -- (4.2,0.4) -- (5.2,2.2) -- (1.0,1.8) -- cycle;
\node[lbl, black!60, anchor=west] at (4.5,2.5) {【平面名，如：赤道面/可行域】};

% ---- 坐标轴（原点按图形重心放） ----
\coordinate (O) at (2.6,0.9);
\draw[axis] (O) -- ++(2.2,0.15) node[lbl, below right] {$X$【轴名】};
\draw[axis] (O) -- ++(-1.5,-0.7) node[lbl, below left] {$Y$【轴名】};
\draw[axis] (O) -- ++(0,2.1)  node[lbl, above] {$Z$【轴名】};
\node[lbl, black!70, below=3pt of O, xshift=-2mm] {原点 $O$};

% ---- 主矢量（如：状态向量/作用方向） ----
\coordinate (P) at ($(O)+(1.5,1.5)$);
\draw[vec] (O) -- (P) node[midway, above left, lbl] {$\mathbf{R}$【矢量名】};
\fill (P) circle (1.2pt);
\node[lbl, right=1pt of P] {【端点含义】};

% ---- 角度标注（angles 库：弧+角度符号） ----
\draw[aux] (O) -- ++(1.9,0.55) coordinate (A1);
\pic[draw=black!60, line width=0.5pt, angle radius=9mm,
     angle eccentricity=1.4, "$\theta$【角名】", lbl] {angle = A1--O--P};

% ---- 辅助虚线（投影/参考线） ----
\draw[aux] (P) -- ($(O)+(1.5,0)$) node[lbl, midway, right, black!60] {投影};

% ---- 参数说明框（图内统计/定义，标题留给论文 caption） ----
\node[callout, anchor=north west] at (-0.3,3.4)
  {$a$：【参数 1 含义】\quad $e$：【参数 2 含义】\\
   $\omega$：【参数 3 含义】};

\end{tikzpicture}
\end{document}
